\section{Typography and its applications}
\subsection{Theoretical Foundations}
Typography operates as a cognitive and semiotic system shaping how users perceive and process content. Functionally, typefaces influence readability and cognitive load in digital contexts, while symbolically, typographic form conveys meaning, authority, and brand identity before linguistic comprehension occurs \cite{lidwell2010, bringhurst2013}.
\subsection{Applied Technical Systems in Web Interface Design}

In practical web interface design, typography functions as both a visual system and a performance-sensitive technical component, directly affecting readability, loading speed, and user perception across devices.

\begin{itemize}

    \item \textbf{Variable Font Architecture ($\text{WOFF2}$):} Variable fonts consolidate multiple typographic styles (regular, bold, condensed) into a single binary file with continuous axes such as Weight ($wght$), Width ($wdth$), and Slant ($slnt$). In real-world applications, this reduces font payload size, minimizes HTTP requests, and allows designers to fine-tune typographic emphasis dynamically (increasing weight for accessibility or emphasis states) using \texttt{font-variation-settings}, thereby improving Core Web Vitals and perceived interface responsiveness \cite{mdnVariableFonts}.

    \item \textbf{Fluid Typography and Viewport-Based Scaling:} Rather than relying on fixed breakpoints, modern interfaces employ mathematically defined scaling functions to ensure consistent typographic hierarchy across screen sizes. The CSS function

    \begin{equation}
        \text{font-size} = \text{clamp}(f_{min}, V_{scaling}, f_{max})
    \end{equation}

    enables text to scale smoothly with the viewport while remaining bounded for legibility, ensuring that headings, body text, and UI labels maintain proportional relationships on both small mobile screens and large desktop displays.

\end{itemize}


\subsection{Critical Analysis: Generative Trends and Information Equity}
Typography is transitioning toward generative, data-responsive systems enabled by algorithmic methods such as L-systems, which require a careful balance between expressive flexibility and functional readability, while emphasizing ethical responsibility in ensuring information equity through accessible typographic standards.
